\scanpage{61}
\mspage{319}

Considérons maintenant l'hom \unsure{(3)}
\[
\tag{18}
\left\{
\begin{aligned}
\widehat\varphi:\quad &\widehat\Pi_{0,3}\longrightarrow
  \widehat{\mathfrak S}_{0,3}^{+}\\
&\ell_0\longmapsto\ell'_0=\sigma_\infty\rho=\varepsilon_0\\
&\ell_1\longrightarrow\ell'_1=\sigma_\infty\\
&\ell_\infty\longmapsto\ell'_\infty=\rho^{-1}
\end{aligned}
\right.
\]
et écrivons qu'il commute, \unsure{en tant} qu'hom. ext., à $u$.
Cela signifie donc qu'il existe un
\[
\tag{19}
\lambda\in\widehat{\mathfrak S}_{0,3}^{+}
\]
tel que l'on ait
\[
\tag{20}
\widehat\varphi(\unsure{ux})=\lambda\,u\widehat\varphi(x)\,\lambda^{-1}
\qquad \forall x\in\widehat\Pi_{0,3},
\]
et il suffit d'exprimer cette condition pour $x$ égal à l'un des deux
gén. $\ell_0,\ell_1$ de $\widehat\Pi_{0,3}$

Pour $x=\ell_0$ la relation (20) s'écrit
\[
\tag{21}
u(\ell'_0)=\operatorname{int}(\lambda^{-1}\widehat\varphi(g_0))(\ell_0'{}^{p})
\]
ce qui, comparé avec (17), donne
\[
\operatorname{int}(g_0^{-1}\lambda^{-1}\widehat\varphi(g_0))(\ell_0'{}^{p})
=\ell_0'{}^{p}\quad\text{donc}
\]
\[
\deleted{(22)}\qquad
g_0^{-1}\lambda^{-1}\widehat\varphi(g_0)=\ell_0'{}^{\beta}
\qquad(\beta\in\widehat{\mathbb Z})
\]
ou encore
\[
\tag{22}
\boxed{\lambda=\widehat\varphi(g_0)\ell_0'{}^{-\beta}g_0^{-1}}
\inserted{=\widehat\varphi(g_0h_0^{-1})=\widehat\varphi(\sigma_\infty(g_0))}
\]
\margin{$\searrow$ \unsure{cf} plus bas (24 bis)}

Pour $x=\ell_1$ la relation (20) s'écrit

\scanpage{62}

\[
\tag{23}
u(\ell'_1)=\operatorname{int}(\lambda^{-1}\widehat\varphi(g_1))(\ell_1'{}^{p})
\]
\margin{où $(g_1=\sigma_\infty(g_0)=\sigma_\infty g_0\sigma_\infty\unsure{)}$}
ou encore
\[
\sigma_\infty=\operatorname{int}(\lambda^{-1}\widehat\varphi(g_1))(\sigma_\infty)
\]
\margin{Car $\sigma_\infty^{p}=\sigma_\infty$ car
\inserted{$\sigma_\infty^2=1$ et} \unsure{l'image} de $p$ dans
$(\mathbb Z/2\mathbb Z)^*$ est $1$ !}
ou encore $\sigma_\infty$ commute à
\[
\lambda^{-1}\widehat\varphi(g_1)
\overset{\inserted{\text{cf (22)}}}{=}
g_0\ell_0'{}^{\beta}\widehat\varphi
 \bigl(\underbrace{g_0^{-1}\sigma_\infty(g_0)}_{h_0^{-1}}\bigr)
=g_0\ell_0'{}^{\beta}\widehat\varphi(h_0^{-1})
\]
ou encore à son inverse $\widehat\varphi(h_0)\ell_0'{}^{-\beta}g_0^{-1}$
ce qui s'écrit aussi
\[
\underbrace{\sigma_\infty
 \bigl(\widehat\varphi(h_0).\ell_0'{}^{-\beta}.g_0^{-1}\bigr)}
=\widehat\varphi(h_0).\ell_0'{}^{-\beta}.g_0^{-1}
\]
\[
\sigma_\infty(\widehat\varphi(h_0)).
\underbrace{\varepsilon_1^{-\beta}}_{(\rho\sigma_\infty)^{-\beta}}.
\sigma_\infty(g_0)^{-1}
=\widehat\varphi(h_0)\ell_0'{}^{-\beta}g_0^{-1}
\]
\[
\sigma_\infty(\widehat\varphi(h_0))
 (\underbrace{\rho\sigma_\infty}_{\varepsilon_1})^{-\beta}h_0
=\widehat\varphi(h_0).\varepsilon_0^{-\beta}
\]
\[
\tag{24}
\boxed{
\overbrace{\varepsilon_1^{-\beta}h_0\varepsilon_0^{\beta}}
 ^{\inserted{\sigma_\infty(g_0\varepsilon_0^{\beta})^{-1}(g_0\varepsilon_0^{\beta})}}
=\sigma_\infty(\widehat\varphi(h_0))^{-1}\widehat\varphi(h_0)
}
\]
\margin{\unsure{Correct?} $\searrow$}
\margin{\uline{NB.} Si $g_0$ est remplacé par
\inserted{$g'_0=$} $g_0\ell_0^{\lambda}$ $(\lambda\in\widehat{\mathbb Z})$,
$h_0$ est remplacé par $h'_0=\ell_1^{-\lambda}\unsure{h_0|g_0}\ell_0^{\lambda}$,
et on \unsure{trouve} que $\beta$ \unsure{est} \unsure{''} remplacé par
$\beta'=\beta-3\lambda$.
Ceci montre \inserted{?} qu'on peut de \uline{façon unique}
\illegible{1} \uline{normaliser} le choix de $g_0$, de telle façon que l'on ait
$\beta\in\{0,1,2\}$.
(Il faudrait élucider si $\beta$ prend effectivement toutes les
\unsure{trois} valeurs mod $3$\ldots)}
\margin{\uuline{NB} Cette équation équivaut à l'équation
\uline{plus simple}
\[
\text{(24 bis)}\qquad
\boxed{\widehat\varphi(h_0)=\sigma_\infty^{\unsure{\varepsilon}}g_0\varepsilon_0^{\beta}}
\]}
\[
\unsure{(25)}\qquad
\left\{
\begin{aligned}
&h_0\text{ donné par (8) en fonction de }g_0\\
&\varepsilon_0=\ell'_0=\sigma_\infty\rho
   &&(\varepsilon_0^2=\ell_0)\\
&\varepsilon_1=\rho\sigma_\infty=\sigma_0\rho
   &&(\varepsilon_1^2=\ell_1)
\end{aligned}
\right.
\]

Posant, pour $\deleted{\unsure{g}}\ \mu\in
\widehat{\mathfrak S}_{0,3}^{+}$,
\[
\tag{26}
\widetilde\mu=\sigma_\infty(\mu)^{-1}\mu
=\sigma_\infty\mu^{-1}\sigma_\infty\mu=[\sigma_\infty,\mu]
\]
(de sorte que $h_0=\widetilde g_0$) la formule \unsure{(24)|(25)} prend la forme
\[
\tag{27}
\varepsilon_1^{-\beta}\unsure{\widetilde g_0}\varepsilon_0^\beta
=\bigl(\widehat\varphi(\widetilde g_0)\bigr)^{\sim}.
\]

\scanpage{63}
\mspage{320}

En résumé :

\uline{Proposition} Considérons un autom ext. $u$ du groupe à lacets
$\widehat\Pi_{0,3}$, réalisé par $g_0\in\widehat\Pi_{0,3}$
\inserted{mod $L_0$,}
et $p\in\widehat{\mathbb Z}^{*}$,
(et $\gamma\in\widehat{\mathbb Z}^{*}$ uniquement déterminé
\unsure{par termes|en termes} de $g_0$ et $p$)
par les formules (11), au moyen d'un autom effectif
\inserted{(noté encore $u$)} qui commute à $\sigma_\infty$ ---
les données $g_0,p$ (et $\gamma$) \unsure{étant} \inserted{donc}
assujetties à vérifier les conditions (9) et (10) ci-dessus,
plus la condition d'engendrement (11 bis).
\inserted{(On prolonge $u$ \unsure{en} autom de
$\widehat{\mathfrak S}_{0,3}^{+}$
par $\operatorname{int}(u)$, cf (15))}
Pour que l'hom extérieur
$\widehat\varphi:\widehat\Pi_{0,3}\longrightarrow
\widehat{\mathfrak S}_{0,3}^{+}$
défini par (18) commute à l'action ext. de $u$,
il f. et s. que l'on puisse trouver un $\beta\in\widehat{\mathbb Z}$,
tel qu'on ait (24), ou ce qui revient au $\overline{\mathrm m}$, (27),
\margin{\unsure{mieux :} \uuline{24 bis}}

Cette condition \inserted{(pour $\beta$ fixé)} équivaut aussi à la formule
(20) pour $x=\ell_\infty$ ($\lambda$ \unsure{étant} tjrs donné par (22)),
\deleted{\unsure{(2}} i.e. à la formule
\[
\tag{28}
u(\ell'_\infty)=\operatorname{int}(\lambda^{-1}\widehat\varphi(g_\infty))
 (\ell_\infty'{}^p),\qquad\text{où}
\]
\[
\tag{29}
g_\infty=g_0\ell_0^{-\gamma}\rho^{-1}(h_0)
\qquad(\text{et }\lambda\text{ donné par (22)}).
\]

\scanpage{64}

formule qui précise donc la formule \unsure{présomptive} (18). Elle s'écrit aussi, plus explicitement
\[
\rho^{-1}g^{-1}
=\operatorname{int}\bigl(g_0\varepsilon_0^{\beta}
\widehat\varphi(\ell_0^{\gamma}\rho^{-1}(h_0))\bigr)
(\rho^{-p})
\]
\inserted{\[
g_0\varepsilon_0^{\beta+\gamma}\widehat\varphi(\rho^{-1}(h_0))
\]}
i.e. (en prenant les inverses, et explicitant $g$)
\[
(30)\qquad
\boxed{
\deleted{\illegible{1}}\inserted{h_0^{-1}\ell_1^{\gamma}\rho}
=\operatorname{int}\bigl(
\deleted{\unsure{g_0}}\varepsilon_0^{\deleted{\unsure{\beta+\gamma}}\inserted{\alpha}}
\widehat\varphi(\rho^{-1}(h_0))\bigr)
(\rho^{p\deleted{\unsure{-p}}})
}
\]
\margin{\uline{NB} On pose $\alpha=\beta+\gamma$. L'existence \unsure{de} $\beta\in\widehat{\mathbb Z}$ vérifiant (24), équivaut à celle \unsure{de} $\alpha\in\widehat{\mathbb Z}$ vérifiant (30).}

L'équivalence de (24) et (30) peut paraître un peu étrange, du fait \inserted{p.ex.} que dans (24) l'exposant $p$ n'intervient pas explicitement, mais seulement $h_0$, tandis que la formule (30) fait intervenir la valeur de $p\in\widehat{\mathbb Z}^{*}$ mod $3$ (qui donne $\rho^p=\rho$ ou $\rho^p=\rho^{-1}$). Mais en tout état de cause, que ce soit sous la forme (24) \unsure{où|ou} sous la forme (30), il semble bien que la condition en question soit hautement non tautologique !

Pour bien faire, il faudrait \unsure{d'abord} expliciter les conditions congruentielles du \S~39 en fonction de $g_0,p$.

\scanpage{65}
\mspage{321}

Je voudrais \unsure{comparer} les autom. ext. de $\widehat\Pi_{0,3}$, et ceux de $\widehat{\mathfrak S}_{0,3}^{+}$, tenant compte de la suite exacte
\[
(31)\qquad
\diagram{\begin{tikzcd}[ampersand replacement=\&]
1\arrow[r,no head]\&\widehat\Pi_{0,3}\arrow[r,no head]\&
\widehat{\mathfrak S}_{0,3}^{+}\arrow[r]\arrow[d,draw=none,"{\inserted{\unsure{\wr}}}" description]\&
\mathfrak S_{3}\arrow[r,no head]\&1\\
\&\&\inserted{\widehat{\Pi'}_{0,3}}
\end{tikzcd}}
\]

\deleted{\uline{Proposition}} \inserted{Notons} que tt autom. ext. \inserted{$u$} de $\widehat{\mathfrak S}_{0,3}^{+}$ qui \unsure{invarie} le sous-groupe distingué $\widehat\Pi_{0,3}$, et induit dans le quotient $\mathfrak S_{3}$ un autom. extérieur trivial, peut s'écrire comme déduit d'un autom. \inserted{$u_0$} de $\widehat{\mathfrak S}_{0,3}^{+}$, \unsure{stabilisant} $\widehat\Pi_{0,3}$, et induisant l'\uline{identité} sur le quotient $\mathfrak S_{3}$ --- et $u_0$ est alors déterminé mod composition avec un élément $\operatorname{int}(g)$, où $g\in\widehat{\mathfrak S}_{0,3}^{+}$ \unsure{est} tel que son image dans $\mathfrak S_{3}$ soit central, i.e. $=1$, i.e. $g\in\widehat\Pi_{0,3}$. Donc
\[
(32)\qquad
\begin{gathered}
\text{groupe des automorphismes}\\
\text{ext. de $\widehat{\mathfrak S}_{0,3}^{+}$ qui}\\
\text{induisent dans $\mathfrak S_{3}$ un autom.}\\
\text{ext. trivial (donc \unsure{invarient}}\\
\widehat\Pi_{0,3})
\end{gathered}
\simeq
\begin{gathered}
\text{groupe des}\\
\text{autom. de}\\
\widehat{\mathfrak S}_{0,3}^{+}\text{ qui}\\
\text{induisent}\\
\text{l'identité dans}\\
\text{son quotient $\mathfrak S_{3}$,}\\
\text{\uline{mod} le sous-}\\
\text{groupe des}\\
\text{autom. int.}\\
\text{par des $g\in\widehat\Pi_{0,3}$}
\end{gathered}
\]

\scanpage{66}

Un \inserted{tel} autom $u_0$ est-il connu par sa restriction à $\widehat\Pi_{0,3}$ ? En d'autres termes, un autom. de $\widehat{\mathfrak S}_{0,3}^{+}$ \deleted{qui} ($\simeq\mathrm{Sl}(\unsure{2|1},\mathbb Z)^{\unsure{\prime}\wedge}$) qui induit l'identité sur $\widehat\Pi_{0,3}$ et sur $\mathfrak S_{3}$, est-il trivial ? Il suffit de prouver que sa restriction à l'image inverse \inserted{$\widehat\Pi_{\unsure{0,\infty|0,3}}\cdot\{1,\unsure{\sigma_\infty}\}$} de $\{1,\dot\sigma_\infty\}\subset\mathfrak S_{3}$ est \unsure{trivial|triviale} (par symétrie ce sera aussi le cas en remplaçant $\dot\sigma_\infty$ par $\dot\sigma_0,\dot\sigma_1$, or ces trois sous groupes engendrent $\widehat{\mathfrak S}_{0,3}^{+}$. On aura $f(\sigma_\infty)=\sigma_\infty g$, avec $g\in\widehat\Pi_{0,3}$ [et on doit avoir $\sigma_\infty g\sigma_\infty g=1$ i.e.
\[
g\cdot\sigma_\infty(g)=1
\]
]
de plus \deleted{pour $x\in\Pi_{0,3}$ \unsure{on a}} \inserted{on devra avoir que}
\[
\operatorname{int}(\sigma_\infty g)\vert\widehat\Pi_{0,3}
=\deleted{\unsure{f(x\sigma_\infty)}}\operatorname{int}(\sigma_\infty)\vert\widehat\Pi_{0,3}
\]
d'où \deleted{$g\in$}
\[
\operatorname{int}(g)\vert\widehat\Pi_{0,3}=\mathrm{id}
\quad\text{i.e.}\quad
g\in\unsure{\operatorname{Centre}}(\widehat\Pi_{0,3})
\quad\text{i.e.}\quad g=1,
\]
OK.

Mais à quelle condition un autom. de $\widehat\Pi_{0,3}$ \deleted{peut} se prolonge-t-il en un autom. de $\widehat{\mathfrak S}_{0,3}^{+}$ qui \deleted{soit}\inserted{\unsure{induise}} l'identité sur $\mathfrak S_{3}$ ? On voit aussitôt qu'il \uline{faut} pour ceci que l'autom. ext. de $\widehat\Pi_{0,3}$ qu'il induit commute à l'action de $\mathfrak S_{3}$.

\scanpage{67}
\mspage{\unsure{322|332}}

C'est suffisant, car soit $\Gamma\inserted{{}\subset\mathrm{Aut\,ext}(\widehat{\Pi}_{0,3})}$ le groupe des autom.
ext. de $\widehat{\Pi}_{0,3}$ qui ont cette propriété (on ne
se préoccupe pas pour le moment de la
structure à lacets de $\widehat{\Pi}_{0,3}$), on a
$\Gamma\cap\mathfrak S_3=(1)$ (en identifiant $\mathfrak S_3$ à un
\unsure{ss-groupe} de $\mathrm{Aut\,ext}\,\widehat{\Pi}_{0,3}$) car \deleted{\unsure{\mathfrak S}} le centre
de $\mathfrak S_3$ est trivial, donc
\margin{ceci est \unsure{évident|évidemment}}
\[
\deleted{\unsure{\Gamma\times}}\ \mathfrak S_3\times\Gamma\hookrightarrow\mathrm{Autext}(\widehat{\Pi}_{0,3}),
\]
et on a donc une structure d'extension
\[
1\unsure{\longrightarrow}\widehat{\Pi}_{0,3}\longrightarrow\mathcal E\longrightarrow\mathfrak S_3\times\Gamma\longrightarrow1
\tag{33}
\]
(car le centre de $\widehat{\Pi}_{0,3}$ est trivial), ou encore
une structure d'extension
\[
1\longrightarrow\widehat{\mathfrak S}_{0,3}^{+}\unsure{\longrightarrow}\mathcal E\longrightarrow\Gamma\longrightarrow1
\tag{34}
\]
qui définit une action extérieure de $\Gamma$
sur $\widehat{\mathfrak S}_{0,3}^{+}$ --- qui est l'action ext.
cherchée. Ainsi

\uline{Prop} le groupe des autom. ext. de $\widehat{\Pi}_{0,3}$ qui
\unsure{commutent} à $\mathfrak S_3$ est isom. can. au groupe
des autom. ext. de $\widehat{\mathfrak S}_{0,3}^{+}$ qui induisent
dans $\mathfrak S_3$ l'autom. ext. trivial.
\margin{Mais \uline{tout} autom. de $\mathfrak S_3$ est intérieur !}

\scanpage{68}

\uline{Corollaire} le groupe des autom. de $\widehat{\Pi}_{0,3}$ qui
commutent extérieurement à $\mathfrak S_3$, est isom.
au groupe des autom. de $\widehat{\mathfrak S}_{0,3}^{+}$ qui
induisent sur $\mathfrak S_3$ l'autom. triviale (par
restriction de ces derniers)

\uuline{Corollaire}
\[
\diagram{\begin{tikzcd}[ampersand replacement=\&,column sep=small,row sep=large]
1 \arrow[r,no head]
\& \widehat{\Pi}_{0,3} \arrow[r]
\& \mathrm{Aut}\left(\widehat{\Pi}_{0,3}\,\middle|\,
  \begin{array}{l}\mathrm{comm.}\\\mathrm{ext.}^{\unsure{t}}\\\text{à }\mathfrak S_3\end{array}\right) \arrow[r]
\& \mathrm{Autext}\left(\widehat{\Pi}_{0,3}\,\middle|\,
  \begin{array}{l}\mathrm{comm.}\\\text{à }\mathfrak S_3\end{array}\right) \arrow[r]
\& 1
\\
1 \arrow[r,no head]
\& \widehat{\Pi}_{0,3} \arrow[r,no head] \arrow[u,phantom,"{\text{''}}" description]
\& \underbrace{\mathrm{Aut}\left(\widehat{\mathfrak S}_{0,3}^{+}\,\middle|\,
  \begin{array}{l}\text{id. sur le}\\\text{quotient }\mathfrak S_3\end{array}\right)}_{\mathcal E^{!}}
  \arrow[r] \arrow[u,"{\sim\ \text{\unsure{restr.} à }\widehat{\Pi}_{0,3}}"]
\& \mathrm{Autext}\left(\widehat{\mathfrak S}_{0,3}^{+}\,\middle|\,
  \begin{array}{l}\text{\deleted{\unsure{finie}}\inserted{\unsure{trivial}}}\\\text{sur le}\\\text{quotient }\mathfrak S_3\end{array}\right)
  \arrow[r] \arrow[u,"\sim"]
\& \unsure{1}
\end{tikzcd}}
\tag{35}
\]

\deleted{\uline{Corollaire}
Soit $u\in\mathrm{Aut}\deleted{\mathrm{ext}}(\widehat{\mathfrak S}_{0,3}^{+})$, \unsure{induisant}
\unsure{Soit} \illegible{1} $u\in\mathrm{Aut}(\widehat{\mathfrak S}_{0,3}^{+})$ et $\unsure{u_0}\unsure{=|:=}u|\widehat{\Pi}_{0,3}$
\margin{dans $\mathcal E$}
l'autom. \inserted{ext.} trivial sur $\mathfrak S_3$. On a donc
\unsure{commutant} \illegible{2}
\[
u(\sigma_\infty)=g\sigma_\infty\quad\text{avec}\quad g\in\widehat{\Pi}_{0,3}\quad\text{\unsure{bien} déterminé},
\]
on a \illegible{1} que $\sigma_\infty$
et $(g\sigma_\infty)^2=1$ i.e. $g\sigma_\infty g\sigma_\infty=1$ i.e. $g\sigma_\infty g$.}

Cette extension $\mathcal E^{!}$ de $\Gamma$ par $\widehat{\Pi}_{0,3}$ est
scindée\inserted{\unsure{non canoniquement}}, le sous-groupe \unsure{splittant|splitting} est formé
des autom. de $\widehat{\Pi}_{0,3}$ qui commutent à $\sigma_\infty$,
ou encore les autom. de $\widehat{\mathfrak S}_{0,3}^{+}$ qui
\deleted{\unsure{commutent}}\inserted{fixent} \unsure{à} $\sigma_\infty$ (et \inserted{qui fixent} $\rho\ \mathrm{mod}\ \widehat{\Pi}_{0,3}$)
\margin{provient du fait qu'il \deleted{\illegible{2}} \inserted{\unsure{existe}} \deleted{\illegible{1}} \unsure{un} \illegible{2} \inserted{$\widehat{\Pi}_{0,3}$-\unsure{conj.}} \illegible{1} d'ordre $2$ au dessus de $\sigma_\infty$ \ldots}
\[
\left\{
\begin{aligned}
u(\sigma_\infty)&=\sigma_\infty\\
u(\rho)&=g.\rho
\end{aligned}
\right.
\qquad(g\in\widehat{\Pi}_{0,3})
\tag{36}
\]

\scanpage{69}
\mspage{323}

\unsure{écrivant} que $u$ est un \unsure{endom.} (car on \unsure{se rappelle}
que $\widehat{\mathfrak S}_{0,3}^{+}$ est défini par les gén. $\sigma_\infty,\rho$
avec les relations
\begin{equation}
\sigma_\infty^2=\rho^3=1
\tag{37}
\end{equation}
\unsure{on trouve} la condition
\begin{equation}
(g\rho)^3=1\qquad\text{i.e.}\qquad g\,\rho(g)\,\rho^2(g)=1.
\tag{38}
\end{equation}

\margin{de plus il \unsure{faut} pour que $\sigma_\infty$ et $g\rho$
\unsure{engendrent} $\unsure{\widehat\Pi_{0,3}}$, avec seules relations
$\sigma_\infty^{\unsure{2}}=(g\rho)^3=1$.}

\uline{Corollaire}. Soit $u\in\mathcal E^{!}$ un autom. de
$\widehat{\mathfrak S}_{0,3}^{+}$ qui induit l'identité sur $\mathfrak S_3$.
Conditions équivalentes:

a) $\exists\ p\in\widehat{\mathbb Z}^{*}$ tel que
$u(\underset{\varepsilon_0}{\uline{\ell'_0}})$ conjugué dans
$\widehat{\mathfrak S}_{0,3}^{+}$ à $\ell_0'{}^p$
(ce $p$ est alors unique\inserted{, et on peut conjuguer par un \unsure{élément} de $\widehat\Pi_{0,3}$}).

b) $\exists\ p\in\widehat{\mathbb Z}^{*}$ tel que $u(\ell_0)$ soit conjugué dans
$\widehat\Pi_{0,3}$ à $\ell_0^p$ (ce $p$ est alors unique).

c) $u_0\overset{\mathrm{dfn}}{=}u\mid\widehat\Pi_{0,3}$ respecte la structure à lacets de $\widehat\Pi_{0,3}$.

De plus, le $p$ dans a) et b) est le $\overline{\mathrm m}$.

\uline{\unsure{Dém}}. Si $u(\ell'_0)=\operatorname{int}(g_0)(\ell_0'{}^p)$, avec
$g_0\in\widehat{\mathfrak S}_{0,3}^{+}$,
élevant au carré on \unsure{aura} $u(\ell_0)=\operatorname{int}(g_0)(\ell_0^p)$,
cela montre que l'image $\dot g_0$ de $g_0$ dans $\mathfrak S_3$
\unsure{invarie} $0$ i.e. est \deleted{\illegible{1}} $1$ ou $\dot\sigma_0$,
\deleted{donc} si c'est $1$ on a bien b), sinon on \unsure{met|\deleted{sait q}}
$g_0=h\sigma_0$ \inserted{($h\in\widehat\Pi_{0,3}$)}
\deleted{\unsure{d'où}}
\[
u(\ell_0)=\operatorname{int}(h)\,
\underset{(\sigma_0(\ell_0))^p}{\uline{\sigma_0(\ell_0^p)}}
\]
or on sait que $\sigma_0(\ell_0)=\operatorname{int}(a)\ell_0$
\inserted{(avec $a\in\widehat\Pi_{0,3}$)},
\inserted{d'où $\sigma_0(\ell_0)^p=\operatorname{int}(a)\ell_0^p$},
donc $u(\ell_0)=\operatorname{int}(ha)(\ell_0^p)$, d'où b) avec le $\overline{\mathrm m}$ $p$.

\scanpage{70}

Comme on sait l'unicité de $p$ dans b), dans a) aussi. Si on a
\begin{equation}
u_0(\ell_0)=\operatorname{int}(g_0)(\ell_0^p)
\qquad (g_0\in\widehat\Pi_{0,3})
\tag{39}
\end{equation}
on en déduit, écrivant
\[
\deleted{\rho u_0\rho^{-1}(\ell_0)\unsure{=}}
\qquad
\rho u_0\rho^{-1}=\operatorname{int}(g^{-1})u_0
\qquad (\unsure{a|g}\in\widehat\Pi_{0,3})
\]
\[
\deleted{u_0 \ell_1}
\]
\[
u_0(\ell_1)=\operatorname{int}(\unsure{a|g})
\underbrace{\rho\underbrace{u_0\underbrace{\rho^{-1}(\ell_1)}_{\ell_0}}_{
\operatorname{int}(g_0)(\ell_0^p)}}_{
\operatorname{int}(\rho(g_0))(\ell_1^p)}
=\operatorname{int}(g\rho(g_0))(\ell_1^p)
\]
i.e.
\begin{equation}
u_0(\ell_1)=\operatorname{int}(g_1)(\ell_1^p)
\qquad\text{avec}\quad g_1=g\rho(g_0)
\tag{40}
\end{equation}
et de $\overline{\mathrm m}$
\begin{equation}
\unsure{u_0|u}(\ell_\infty)=\operatorname{int}(g_\infty)(\ell_\infty^p)
\qquad\text{avec}\quad
\begin{aligned}
g_\infty&=g\rho(g_1)\\
&=g\rho(g)\rho^2(g_0)
\end{aligned}
\tag{41}
\end{equation}
donc b) $\Rightarrow$ c). Enfin \unsure{prouvons} c) $\Rightarrow$ a).

On a par hypothèse
\[
\begin{array}{rcl}
u_0(\ell_0)&=&\operatorname{int}(g)\ell_0^p\\
\text{''}&&\text{''}\\
u(\varepsilon_0)^2&&
\operatorname{int}(g)(\unsure{\ell|\varepsilon}_0^{\unsure{2p|p}})^2
=(\operatorname{int}g(\unsure{\ell_0^p|\varepsilon_0^p}))^2
\end{array}
\]
je dis qu'on a alors
\[
u(\varepsilon_0)=\operatorname{int}(g)(\varepsilon_0^p).
\]
Il suffit de \deleted{voir}\inserted{\unsure{se rappeler}} que
\deleted{\unsure{tt}}\inserted{le seul} \unsure{élément} de
$\widehat{\mathfrak S}_{0,3}^{+}$ de carré égal à $\ell_0$ est $\varepsilon_0$\ldots\quad OK

\margin{\unsure{vérifier}}
